% --- Archon physics formalization source begin ---
% archon:physics
% archon:covers PhyXMiniProblems/problem_phyx_mini_0767.lean
% archon:source-report reports/phyx_mini/problem_phyx_mini_0767.source.json

\chapter{Physics problem phyx\_mini\_0767}
\label{ch:PhyXMiniProblems_problem_phyx_mini_0767}

\paragraph{Problem source.}
\#\# Physical scenario

Two identical masses are released from rest in a smooth hemispherical bowl of radius \textbackslash{}(R\textbackslash{}) from the positions shown in figure. Ignore friction between the masses and the surface of the bowl.

\#\# Question

If the masses stick together when they collide, how high above the bottom of the bowl will they go after colliding?

\#\# Answer choices

- A: A: R/3

- B: B: R/4

- C: C: R/2

- D: D: R/5

\#\# Figure caption (auxiliary; use the image as primary evidence)

The image depicts a semicircular setup with two squares and a line. Here's a detailed breakdown:

- **Semicircle**: The image shows a semicircle with a thickened gray border.
- **Squares**: 
  - One square is positioned near the left edge of the semicircle, colored orange with a darker border.
  - Another square is located at the bottom of the semicircle, similarly colored and bordered.
- **Line**: A line extends from the center of the semicircle towards the right side, ending outside the semicircle. 
- **Text**: The letter "R" is annotated next to the line, with an arrow indicating the direction from the center to the edge.

The layout suggests a geometric or physical scenario, possibly related to forces or distances.

\#\# Dataset metadata

Momentum and Collisions; Physical Model Grounding Reasoning, Multi-Formula Reasoning

\paragraph{Recorded answer/context.}
B: B: R/4

\paragraph{Figure/image path.}
/root/proposal\_for\_physic/science-mango/phyx\_mini\_run/phyx\_data/test\_image/767.png

\paragraph{Formalization target.}
create a compiling Lean file with sorry bodies at `PhyXMiniProblems/problem\_phyx\_mini\_0767.lean`. The Lean declarations must preserve the physical quantities, dimensions or dimensional roles, figure labels, governing-law hypotheses, and final relation expressed by this problem.
Use Mathlib/Physlib names found through LeanExplore where available. If a physics API is missing, introduce faithful local abstractions rather than scalar placeholder aliases.

\begin{theorem}[Physics formalization target]
\label{thm:physics:phyx_mini_0767:target}
\lean{PhyXMiniProblems.ProblemPhyXMini0767.problem_phyx_mini_0767}
\uses{def:physics:phyx-mini-0767:phyxminiproblems-problemphyxmini0767-lengthreadout, def:physics:phyx-mini-0767:phyxminiproblems-problemphyxmini0767-bowlcollisionsetup, def:physics:phyx-mini-0767:phyxminiproblems-problemphyxmini0767-matchesbowlcollisionscenario, def:physics:phyx-mini-0767:phyxminiproblems-problemphyxmini0767-matchessuppliedfigure, def:physics:phyx-mini-0767:phyxminiproblems-problemphyxmini0767-matchesinitialandcollisiongeometry, def:physics:phyx-mini-0767:phyxminiproblems-problemphyxmini0767-hasphysicalbowlparameters, def:physics:phyx-mini-0767:phyxminiproblems-problemphyxmini0767-satisfiesbowlcollisionconservationlaws, def:physics:phyx-mini-0767:phyxminiproblems-problemphyxmini0767-iscorrectanswer, def:physics:phyx-mini-0767:phyxminiproblems-problemphyxmini0767-recordedanswerchoice}
The assigned autoformalize agent should translate this physics problem into Lean declarations in the covered file, with theorem and lemma proof bodies written as `by sorry`.
\end{theorem}
\begin{proof}
This is an autoformalization task, not a proof task. Produce statements that can later be proved by the physics prover without weakening the physical model.
\end{proof}

% --- Archon declaration topology begin ---
\section{Declaration topology}
\begin{definition}[Length Quantity]
  \label{def:physics:phyx-mini-0767:phyxminiproblems-problemphyxmini0767-lengthquantity}
  \lean{PhyXMiniProblems.ProblemPhyXMini0767.LengthQuantity}
  A nonnegative physical length, independent of the unit used to read it
\end{definition}
\begin{proof}
  This is the declared model definition of Length Quantity.
\end{proof}
\begin{definition}[Mass Quantity]
  \label{def:physics:phyx-mini-0767:phyxminiproblems-problemphyxmini0767-massquantity}
  \lean{PhyXMiniProblems.ProblemPhyXMini0767.MassQuantity}
  A nonnegative physical mass, independent of the unit used to read it
\end{definition}
\begin{proof}
  This is the declared model definition of Mass Quantity.
\end{proof}
\begin{definition}[Acceleration Magnitude Quantity]
  \label{def:physics:phyx-mini-0767:phyxminiproblems-problemphyxmini0767-accelerationmagnitudequantity}
  \lean{PhyXMiniProblems.ProblemPhyXMini0767.AccelerationMagnitudeQuantity}
  A nonnegative acceleration magnitude with physical dimension L T⁻²
\end{definition}
\begin{proof}
  This is the declared model definition of Acceleration Magnitude Quantity.
\end{proof}
\begin{definition}[Tangential Velocity Quantity]
  \label{def:physics:phyx-mini-0767:phyxminiproblems-problemphyxmini0767-tangentialvelocityquantity}
  \lean{PhyXMiniProblems.ProblemPhyXMini0767.TangentialVelocityQuantity}
  A signed tangential velocity. The real carrier retains the sign needed by the one-dimensional momentum balance, while WithDim retains dimension L T⁻¹
\end{definition}
\begin{proof}
  This is the declared model definition of Tangential Velocity Quantity.
\end{proof}
\begin{definition}[length Readout]
  \label{def:physics:phyx-mini-0767:phyxminiproblems-problemphyxmini0767-lengthreadout}
  \lean{PhyXMiniProblems.ProblemPhyXMini0767.lengthReadout}
  \uses{def:physics:phyx-mini-0767:phyxminiproblems-problemphyxmini0767-lengthquantity}
  Scalar readout of a physical length in a coherent choice of units
\end{definition}
\begin{proof}
  This is the declared model definition of length Readout.
\end{proof}
\begin{definition}[mass Readout]
  \label{def:physics:phyx-mini-0767:phyxminiproblems-problemphyxmini0767-massreadout}
  \lean{PhyXMiniProblems.ProblemPhyXMini0767.massReadout}
  \uses{def:physics:phyx-mini-0767:phyxminiproblems-problemphyxmini0767-massquantity}
  Scalar readout of a physical mass in a coherent choice of units
\end{definition}
\begin{proof}
  This is the declared model definition of mass Readout.
\end{proof}
\begin{definition}[acceleration Readout]
  \label{def:physics:phyx-mini-0767:phyxminiproblems-problemphyxmini0767-accelerationreadout}
  \lean{PhyXMiniProblems.ProblemPhyXMini0767.accelerationReadout}
  \uses{def:physics:phyx-mini-0767:phyxminiproblems-problemphyxmini0767-accelerationmagnitudequantity}
  Scalar readout of an acceleration magnitude in coherent units
\end{definition}
\begin{proof}
  This is the declared model definition of acceleration Readout.
\end{proof}
\begin{definition}[tangential Velocity Readout]
  \label{def:physics:phyx-mini-0767:phyxminiproblems-problemphyxmini0767-tangentialvelocityreadout}
  \lean{PhyXMiniProblems.ProblemPhyXMini0767.tangentialVelocityReadout}
  \uses{def:physics:phyx-mini-0767:phyxminiproblems-problemphyxmini0767-tangentialvelocityquantity}
  Scalar readout of a signed tangential velocity in coherent units
\end{definition}
\begin{proof}
  This is the declared model definition of tangential Velocity Readout.
\end{proof}
\begin{definition}[gravitational Potential Energy Readout]
  \label{def:physics:phyx-mini-0767:phyxminiproblems-problemphyxmini0767-gravitationalpotentialenergyreadout}
  \lean{PhyXMiniProblems.ProblemPhyXMini0767.gravitationalPotentialEnergyReadout}
  \uses{def:physics:phyx-mini-0767:phyxminiproblems-problemphyxmini0767-lengthquantity, def:physics:phyx-mini-0767:phyxminiproblems-problemphyxmini0767-massquantity, def:physics:phyx-mini-0767:phyxminiproblems-problemphyxmini0767-accelerationmagnitudequantity, def:physics:phyx-mini-0767:phyxminiproblems-problemphyxmini0767-lengthreadout, def:physics:phyx-mini-0767:phyxminiproblems-problemphyxmini0767-massreadout, def:physics:phyx-mini-0767:phyxminiproblems-problemphyxmini0767-accelerationreadout}
  The scalar readout of m g h. Its factors are dimension-tagged quantities, so this helper is an energy readout rather than an untyped physical primitive
\end{definition}
\begin{proof}
  This is the declared model definition of gravitational Potential Energy Readout.
\end{proof}
\begin{definition}[translational Kinetic Energy Readout]
  \label{def:physics:phyx-mini-0767:phyxminiproblems-problemphyxmini0767-translationalkineticenergyreadout}
  \lean{PhyXMiniProblems.ProblemPhyXMini0767.translationalKineticEnergyReadout}
  \uses{def:physics:phyx-mini-0767:phyxminiproblems-problemphyxmini0767-massquantity, def:physics:phyx-mini-0767:phyxminiproblems-problemphyxmini0767-tangentialvelocityquantity, def:physics:phyx-mini-0767:phyxminiproblems-problemphyxmini0767-massreadout, def:physics:phyx-mini-0767:phyxminiproblems-problemphyxmini0767-tangentialvelocityreadout}
  The scalar readout of the translational kinetic energy m v² / 2
\end{definition}
\begin{proof}
  This is the declared model definition of translational Kinetic Energy Readout.
\end{proof}
\begin{definition}[tangential Momentum Readout]
  \label{def:physics:phyx-mini-0767:phyxminiproblems-problemphyxmini0767-tangentialmomentumreadout}
  \lean{PhyXMiniProblems.ProblemPhyXMini0767.tangentialMomentumReadout}
  \uses{def:physics:phyx-mini-0767:phyxminiproblems-problemphyxmini0767-massquantity, def:physics:phyx-mini-0767:phyxminiproblems-problemphyxmini0767-tangentialvelocityquantity, def:physics:phyx-mini-0767:phyxminiproblems-problemphyxmini0767-massreadout, def:physics:phyx-mini-0767:phyxminiproblems-problemphyxmini0767-tangentialvelocityreadout}
  The signed one-dimensional tangential momentum readout m v
\end{definition}
\begin{proof}
  This is the declared model definition of tangential Momentum Readout.
\end{proof}
\begin{definition}[Bowl Mass]
  \label{def:physics:phyx-mini-0767:phyxminiproblems-problemphyxmini0767-bowlmass}
  \lean{PhyXMiniProblems.ProblemPhyXMini0767.BowlMass}
  The two identical bodies distinguished by their pictured starting places
\end{definition}
\begin{proof}
  This is the declared model definition of Bowl Mass.
\end{proof}
\begin{definition}[Bowl Location]
  \label{def:physics:phyx-mini-0767:phyxminiproblems-problemphyxmini0767-bowllocation}
  \lean{PhyXMiniProblems.ProblemPhyXMini0767.BowlLocation}
  Distinguished locations in the lower semicircular cross-section
\end{definition}
\begin{proof}
  This is the declared model definition of Bowl Location.
\end{proof}
\begin{definition}[Bowl Geometry]
  \label{def:physics:phyx-mini-0767:phyxminiproblems-problemphyxmini0767-bowlgeometry}
  \lean{PhyXMiniProblems.ProblemPhyXMini0767.BowlGeometry}
  Geometric model of the confining bowl
\end{definition}
\begin{proof}
  This is the declared model definition of Bowl Geometry.
\end{proof}
\begin{definition}[Surface Condition]
  \label{def:physics:phyx-mini-0767:phyxminiproblems-problemphyxmini0767-surfacecondition}
  \lean{PhyXMiniProblems.ProblemPhyXMini0767.SurfaceCondition}
  Contact condition between each mass and the bowl surface
\end{definition}
\begin{proof}
  This is the declared model definition of Surface Condition.
\end{proof}
\begin{definition}[Collision Behavior]
  \label{def:physics:phyx-mini-0767:phyxminiproblems-problemphyxmini0767-collisionbehavior}
  \lean{PhyXMiniProblems.ProblemPhyXMini0767.CollisionBehavior}
  Mechanical behavior of the two bodies at impact
\end{definition}
\begin{proof}
  This is the declared model definition of Collision Behavior.
\end{proof}
\begin{definition}[Figure Label]
  \label{def:physics:phyx-mini-0767:phyxminiproblems-problemphyxmini0767-figurelabel}
  \lean{PhyXMiniProblems.ProblemPhyXMini0767.FigureLabel}
  The literal symbolic label visible beside the radius segment
\end{definition}
\begin{proof}
  This is the declared model definition of Figure Label.
\end{proof}
\begin{definition}[Hemispherical Bowl Figure]
  \label{def:physics:phyx-mini-0767:phyxminiproblems-problemphyxmini0767-hemisphericalbowlfigure}
  \lean{PhyXMiniProblems.ProblemPhyXMini0767.HemisphericalBowlFigure}
  \uses{def:physics:phyx-mini-0767:phyxminiproblems-problemphyxmini0767-bowlmass, def:physics:phyx-mini-0767:phyxminiproblems-problemphyxmini0767-bowllocation, def:physics:phyx-mini-0767:phyxminiproblems-problemphyxmini0767-figurelabel}
  Qualitative information in the primary raster phyx\_data/test\_image/767.png. The image is a lower semicircular cross-section with one square at the left rim, one at the bottom, and a center-to-wall segment marked R
\end{definition}
\begin{proof}
  This is the declared model definition of Hemispherical Bowl Figure.
\end{proof}
\begin{definition}[Bowl Collision Setup]
  \label{def:physics:phyx-mini-0767:phyxminiproblems-problemphyxmini0767-bowlcollisionsetup}
  \lean{PhyXMiniProblems.ProblemPhyXMini0767.BowlCollisionSetup}
  \uses{def:physics:phyx-mini-0767:phyxminiproblems-problemphyxmini0767-lengthquantity, def:physics:phyx-mini-0767:phyxminiproblems-problemphyxmini0767-massquantity, def:physics:phyx-mini-0767:phyxminiproblems-problemphyxmini0767-accelerationmagnitudequantity, def:physics:phyx-mini-0767:phyxminiproblems-problemphyxmini0767-tangentialvelocityquantity, def:physics:phyx-mini-0767:phyxminiproblems-problemphyxmini0767-bowlmass, def:physics:phyx-mini-0767:phyxminiproblems-problemphyxmini0767-bowllocation, def:physics:phyx-mini-0767:phyxminiproblems-problemphyxmini0767-bowlgeometry, def:physics:phyx-mini-0767:phyxminiproblems-problemphyxmini0767-surfacecondition, def:physics:phyx-mini-0767:phyxminiproblems-problemphyxmini0767-collisionbehavior, def:physics:phyx-mini-0767:phyxminiproblems-problemphyxmini0767-hemisphericalbowlfigure}
  Independent physical quantities for the experiment. In particular, maximumRiseHeight is an unconstrained observable; it is not defined from the radius or from an answer choice
\end{definition}
\begin{proof}
  This is the declared model definition of Bowl Collision Setup.
\end{proof}
\begin{definition}[Matches Bowl Collision Scenario]
  \label{def:physics:phyx-mini-0767:phyxminiproblems-problemphyxmini0767-matchesbowlcollisionscenario}
  \lean{PhyXMiniProblems.ProblemPhyXMini0767.MatchesBowlCollisionScenario}
  \uses{def:physics:phyx-mini-0767:phyxminiproblems-problemphyxmini0767-bowlcollisionsetup}
  The qualitative assumptions stated in the problem prose
\end{definition}
\begin{proof}
  This is the declared model definition of Matches Bowl Collision Scenario.
\end{proof}
\begin{definition}[Matches Supplied Figure]
  \label{def:physics:phyx-mini-0767:phyxminiproblems-problemphyxmini0767-matchessuppliedfigure}
  \lean{PhyXMiniProblems.ProblemPhyXMini0767.MatchesSuppliedFigure}
  \uses{def:physics:phyx-mini-0767:phyxminiproblems-problemphyxmini0767-bowlcollisionsetup}
  Facts read directly from the supplied primary figure
\end{definition}
\begin{proof}
  This is the declared model definition of Matches Supplied Figure.
\end{proof}
\begin{definition}[Matches Initial And Collision Geometry]
  \label{def:physics:phyx-mini-0767:phyxminiproblems-problemphyxmini0767-matchesinitialandcollisiongeometry}
  \lean{PhyXMiniProblems.ProblemPhyXMini0767.MatchesInitialAndCollisionGeometry}
  \uses{def:physics:phyx-mini-0767:phyxminiproblems-problemphyxmini0767-lengthreadout, def:physics:phyx-mini-0767:phyxminiproblems-problemphyxmini0767-tangentialvelocityreadout, def:physics:phyx-mini-0767:phyxminiproblems-problemphyxmini0767-bowlcollisionsetup}
  Quantitative geometry and initial conditions supplied by the prose and image. The left-rim height equals the radius, the other mass and the collision are at the zero-height bottom, both bodies are released from rest, and their physical masses are equal. No post-collision speed or rise height is fixed
\end{definition}
\begin{proof}
  This is the declared model definition of Matches Initial And Collision Geometry.
\end{proof}
\begin{definition}[Has Physical Bowl Parameters]
  \label{def:physics:phyx-mini-0767:phyxminiproblems-problemphyxmini0767-hasphysicalbowlparameters}
  \lean{PhyXMiniProblems.ProblemPhyXMini0767.HasPhysicalBowlParameters}
  \uses{def:physics:phyx-mini-0767:phyxminiproblems-problemphyxmini0767-lengthreadout, def:physics:phyx-mini-0767:phyxminiproblems-problemphyxmini0767-massreadout, def:physics:phyx-mini-0767:phyxminiproblems-problemphyxmini0767-accelerationreadout, def:physics:phyx-mini-0767:phyxminiproblems-problemphyxmini0767-tangentialvelocityreadout, def:physics:phyx-mini-0767:phyxminiproblems-problemphyxmini0767-bowlcollisionsetup}
  Positivity and orientation conditions selecting the physical branch
\end{definition}
\begin{proof}
  This is the declared model definition of Has Physical Bowl Parameters.
\end{proof}
\begin{definition}[Satisfies Bowl Collision Conservation Laws]
  \label{def:physics:phyx-mini-0767:phyxminiproblems-problemphyxmini0767-satisfiesbowlcollisionconservationlaws}
  \lean{PhyXMiniProblems.ProblemPhyXMini0767.SatisfiesBowlCollisionConservationLaws}
  \uses{def:physics:phyx-mini-0767:phyxminiproblems-problemphyxmini0767-lengthreadout, def:physics:phyx-mini-0767:phyxminiproblems-problemphyxmini0767-massreadout, def:physics:phyx-mini-0767:phyxminiproblems-problemphyxmini0767-accelerationreadout, def:physics:phyx-mini-0767:phyxminiproblems-problemphyxmini0767-tangentialvelocityreadout, def:physics:phyx-mini-0767:phyxminiproblems-problemphyxmini0767-gravitationalpotentialenergyreadout, def:physics:phyx-mini-0767:phyxminiproblems-problemphyxmini0767-translationalkineticenergyreadout, def:physics:phyx-mini-0767:phyxminiproblems-problemphyxmini0767-tangentialmomentumreadout, def:physics:phyx-mini-0767:phyxminiproblems-problemphyxmini0767-bowlcollisionsetup}
  The generic mechanical laws used in the textbook solution: mechanical energy of the rim mass is conserved during its frictionless descent; the mass initially at the bottom remains at rest until impact; tangential momentum is conserved during the short sticking collision; and mechanical energy of the joined body is conserved from just after impact to its zero-speed turning point. These laws quantify over every coherent unit choice.
\end{definition}
\begin{proof}
  This is the declared model definition of Satisfies Bowl Collision Conservation Laws.
\end{proof}
\begin{lemma}[rim impact velocity squared]
  \label{lem:physics:phyx-mini-0767:phyxminiproblems-problemphyxmini0767-rim-impact-velocity-squared}
  \lean{PhyXMiniProblems.ProblemPhyXMini0767.rim_impact_velocity_squared}
  \uses{def:physics:phyx-mini-0767:phyxminiproblems-problemphyxmini0767-lengthreadout, def:physics:phyx-mini-0767:phyxminiproblems-problemphyxmini0767-accelerationreadout, def:physics:phyx-mini-0767:phyxminiproblems-problemphyxmini0767-tangentialvelocityreadout, def:physics:phyx-mini-0767:phyxminiproblems-problemphyxmini0767-bowlcollisionsetup, def:physics:phyx-mini-0767:phyxminiproblems-problemphyxmini0767-matchesinitialandcollisiongeometry, def:physics:phyx-mini-0767:phyxminiproblems-problemphyxmini0767-hasphysicalbowlparameters, def:physics:phyx-mini-0767:phyxminiproblems-problemphyxmini0767-satisfiesbowlcollisionconservationlaws}
  Frictionless descent from the rim to the bottom gives v\_impact² = 2 g R
\end{lemma}
\begin{proof}
  The conclusion follows from the governing relations and figure readouts named in the statement of rim impact velocity squared.
\end{proof}
\begin{lemma}[common velocity after sticking is half]
  \label{lem:physics:phyx-mini-0767:phyxminiproblems-problemphyxmini0767-common-velocity-after-sticking-is-half}
  \lean{PhyXMiniProblems.ProblemPhyXMini0767.common_velocity_after_sticking_is_half}
  \uses{def:physics:phyx-mini-0767:phyxminiproblems-problemphyxmini0767-tangentialvelocityreadout, def:physics:phyx-mini-0767:phyxminiproblems-problemphyxmini0767-bowlcollisionsetup, def:physics:phyx-mini-0767:phyxminiproblems-problemphyxmini0767-matchesinitialandcollisiongeometry, def:physics:phyx-mini-0767:phyxminiproblems-problemphyxmini0767-hasphysicalbowlparameters, def:physics:phyx-mini-0767:phyxminiproblems-problemphyxmini0767-satisfiesbowlcollisionconservationlaws}
  For two equal masses, with the bottom mass initially stationary, tangential momentum conservation makes the shared post-impact velocity half the incoming rim-mass velocity
\end{lemma}
\begin{proof}
  The conclusion follows from the governing relations and figure readouts named in the statement of common velocity after sticking is half.
\end{proof}
\begin{definition}[Answer Choice]
  \label{def:physics:phyx-mini-0767:phyxminiproblems-problemphyxmini0767-answerchoice}
  \lean{PhyXMiniProblems.ProblemPhyXMini0767.AnswerChoice}
  Labels printed beside the four candidate heights
\end{definition}
\begin{proof}
  This is the declared model definition of Answer Choice.
\end{proof}
\begin{definition}[radius Fraction]
  \label{def:physics:phyx-mini-0767:phyxminiproblems-problemphyxmini0767-answerchoice-radiusfraction}
  \lean{PhyXMiniProblems.ProblemPhyXMini0767.AnswerChoice.radiusFraction}
  \uses{def:physics:phyx-mini-0767:phyxminiproblems-problemphyxmini0767-answerchoice}
  Dimensionless coefficient of the bowl radius printed by each choice
\end{definition}
\begin{proof}
  This is the declared model definition of radius Fraction.
\end{proof}
\begin{definition}[Is Correct Answer]
  \label{def:physics:phyx-mini-0767:phyxminiproblems-problemphyxmini0767-iscorrectanswer}
  \lean{PhyXMiniProblems.ProblemPhyXMini0767.IsCorrectAnswer}
  \uses{def:physics:phyx-mini-0767:phyxminiproblems-problemphyxmini0767-lengthreadout, def:physics:phyx-mini-0767:phyxminiproblems-problemphyxmini0767-bowlcollisionsetup, def:physics:phyx-mini-0767:phyxminiproblems-problemphyxmini0767-answerchoice}
  A choice is correct when its radius fraction equals the physical rise
\end{definition}
\begin{proof}
  This is the declared model definition of Is Correct Answer.
\end{proof}
\begin{definition}[recorded Answer Choice]
  \label{def:physics:phyx-mini-0767:phyxminiproblems-problemphyxmini0767-recordedanswerchoice}
  \lean{PhyXMiniProblems.ProblemPhyXMini0767.recordedAnswerChoice}
  \uses{def:physics:phyx-mini-0767:phyxminiproblems-problemphyxmini0767-answerchoice}
  The answer label recorded in the supplied dataset metadata
\end{definition}
\begin{proof}
  This is the declared model definition of recorded Answer Choice.
\end{proof}
% --- Archon declaration topology end ---
% --- Archon physics formalization source end ---
